% =========================================================
% 1. Abstract
% =========================================================
\section*{一、摘要}
\addcontentsline{toc}{section}{一、摘要 (Abstract)}
在學術人才培育中，研究指導會議不僅推動研究進展，亦是培養學生表達、釐清與修正研究想法的重要互動場域。隨著大型語言模型（Large Language Models, LLMs）逐漸被引入教育與學術輔導情境，已有不少研究嘗試以虛擬系統擔任研究導師或助教。然而，現有多數基於 LLM 的虛擬助教系統仍仰賴通用指導原則或預設角色特性，難以反映研究指導中高度互動且具連續性的實務特性。

本計畫聚焦於現有虛擬研究指導系統在實務應用上的三項關鍵限制：
(一) 研究歷程之邏輯銜接不順：即系統未能有效納入跨會議的研究演進脈絡，包含研究方向的轉折與反覆修正歷程，導致對當前指導目標的理解與實際研究狀態產生落差。
(二) 角色人格漂移（persona drift）：即虛擬助教的指導風格與專業界線容易受到即時語境影響而偏移，進而產生迎合式建議，降低長期互動中回應內容的可靠性 。
(三) 互動臨場感不足：現有系統多採用通用型大型語言模型，回覆傾向標準化且高度中立，難以呈現真實學術情境中個別教授獨有的指導風格與互動習慣，導致對話過程流於機械化，缺乏臨場感。

為回應上述挑戰，本研究提出一套整合多層級記憶協調與角色向量建模之虛擬研究指導對話系統。系統透過分層記憶架構整合即時對話脈絡、近期研究活動摘要與長期研究歷程紀錄，以維持跨會議互動的邏輯一致性；同時，導入即時監控的推理階段介入機制，以在必要時對模型輸出進行干預，將真實研究指導會議中的互動模式融入角色向量之建模，以維持穩定且具臨場感的指導行為。最後，本研究將建立對應長期互動情境的評估方法，針對研究歷程連貫性、指導角色穩定性與互動臨場感等關鍵面向進行系統性驗證，作為評估虛擬研究指導系統是否能貼近真實 mentor--mentee 配對關係的重要依據。

本研究預期發展一套可實作且可評估的虛擬研究指導系統設計架構，使人工智慧能在跨會議、跨研究階段的互動中，維持對研究歷程與指導角色的穩定理解。研究成果可作為未來研究指導與高等教育人機協作系統的設計基礎，使其更適用於實際研究訓練情境，並為相關系統的設計與評估提供具體經驗。此外，透過對長期研究指導會議狀態與互動歷程的分析，本研究亦有助於促進研究指導經驗的系統化整理與分享。